\documentclass[10pt,a4paper]{article}

%% pacotes essenciais e atualizados
\usepackage[T1]{fontenc}      % Essencial para hifenização de palavras acentuadas
\usepackage[utf8]{inputenc}
\usepackage[brazilian]{babel} % Tradução de termos automáticos e regras pt-BR
\usepackage[top=2.5cm, bottom=2.5cm, left=2cm, right=2cm]{geometry} % Controle de margens
\usepackage{graphicx,color,verbatim,Sweave,url,xargs,amsmath,booktabs,longtable,eurosym}

%% novos ambientes requeridos pelo pacote exams
\newenvironment{question}{\item}{}
\newenvironment{solution}{\comment}{\endcomment}
\newenvironment{answerlist}{\renewcommand{\labelenumii}{(\alph{enumii})}\begin{enumerate}}{\end{enumerate}}

%% compatibilidade com pandoc e knitr/rmarkdown
\providecommand{\tightlist}{\setlength{\itemsep}{0pt}\setlength{\parskip}{0pt}}
\providecommand{\pandocbounded}[1]{#1}
\setkeys{Gin}{keepaspectratio}

%% fontes: Helvetica (sem serifa, excelente para provas)
\usepackage{helvet}
\IfFileExists{sfmath.sty}{
  \RequirePackage[helvet]{sfmath}
}{}
\renewcommand{\sfdefault}{phv}
\renewcommand{\rmdefault}{phv}

%% novos comandos para receber variáveis do R
\makeatletter
\newcommand{\Course}[1]{\def\@Course{#1}}
\newcommand{\Semester}[1]{\def\@Semester{#1}}
\newcommand{\ID}[1]{\def\@ID{#1}}

\Course{Disciplina}
\Semester{Ano/Semestre}
\ID{00000}

%% \exinput{header}

\newcommand{\myCourse}{\@Course}
\newcommand{\mySemester}{\@Semester}
\newcommand{\myID}{\@ID}
\makeatother

%% cabeçalhos das páginas secundárias (O nome do aluno aparecerá aqui em todas as páginas)
\markboth{\textnormal{\bf \large \myCourse \: \mySemester \hfill \myID}}%
         {\textnormal{\bf \large \myCourse \: \mySemester \: \myID}}
\pagestyle{myheadings}

\begin{document}

%% página de rosto / cabeçalho principal
\thispagestyle{empty}
{\sf
\begin{center}
    \textbf{\LARGE Universidade Federal de Santa Catarina}\\[0.15cm]
    \textbf{\Large Campus de Curitibanos}\\[0.5cm]
    \textbf{\large Prova de \myCourse \ -- Semestre \mySemester} \hfill \textbf{\myID}
\end{center}

\vspace{0.3cm}

% Quadro exclusivo para assinatura (O nome já foi impresso via ID)
\noindent\fbox{
  \begin{minipage}{\textwidth}
    \vspace{0.3cm}
    \textbf{Assinatura:} \hrulefill \hfill \textbf{Data:} \rule{3.5cm}{0.4pt}
    \vspace{0.3cm}
  \end{minipage}
}
}

\vspace*{0.8cm}

% Início dos exercícios
\begin{enumerate}
\setlength{\itemsep}{1.5em} % Aumenta levemente o espaço entre uma questão e outra
%% \exinput{exercises}
\end{enumerate}

\end{document}
